% title: About the dream of becoming a scientist being a Brazilian
% date: 2023-01
% description: A personal narrative on pursuing science in Brazil.

As most of the people around the world, I did not grow up in a rich family and, to be fair, I also did not grow up in a rich country. The reality of the country where I was born, Brazil, is a Latin American common story: the fight for basic rights and to ensure fair opportunities. But today I will not write about Brazil, but my desire: becoming a scientist being a Brazilian.

When I was born, in 2001, my parents were at the very beginning of their lives being a family: they had recently joined public offices and, then, started to build our home. From their families, that came from the poor countryside, they were the very first to work on stable jobs: my mom as a teacher and my father as policeman.

In 2002, however, something different happened. For the first time, the Brazilian Workers Party (PT) won a presidential election, that made Lula president. Their proposals since the first election after the military dictatorship changed and evolved, but always remained with the same basis, the social transformation.

With this background I lived my childhood. I saw the poverty, but I also saw and lived big changes. I saw my neighborhood streets being paved, my house being improved and my whole family getting better life conditions. I saw my city evolving, homeless people receiving houses and universities being built. Everything seemed to be possible that time.

Since I was a kid, I always loved to study and my mom always encouraged me to keep it as a ``road'' to my life. Therefore, the education opportunities were what I expected the most: I saw simple youngsters like me today getting free opportunities to enroll bachelor courses never before thought to our reality and even going to study abroad. And I could not wait to my time.

In 2013, I lived the most impactful experience of my life. I was selected to represent my state during the IV CNIJMA (a national children's conference for the environment) with a sugarcane burning reduction project to my city. I flew by airplane for the first time, I knew Bras\'ilia, Brazil's capital. More than that, I studied several other subjects, several other projects, with children like me. \emph{The education was beginning to change my life}.

But from what I thought being the beginning of a dream, became the beginning of a dream \textbf{and a fight}. In 2015, Brazil entered a period of political instability, that caused Dilma Rousseff (Brazil's then president) to be impeached. The project that was initiated with Lula since 2003 and then continued with Dilma since 2011, was stopped in 2015, and then, began to be destroyed in 2016.

Going back to my life, in 2016 I got enrolled in the Federal Institute of Paran\'a (IFPR), one of the biggest education projects during Lula's government. My High School was integrated with a Computer Technician course. My teachers were doctors and masters in their areas. The campus had high quality Chemistry, Physics, Informatics and Mechanics labs.

It was during these years at IFPR that I began writing articles and making science and, finally, decided to follow an academic career and fell in love with tech. But it was also during these years that I, again, lived changes: the technical travels, the large projects funding, the structure expansion projects gave place to not even being able to pay the electricity bills due the several reduction of government resources.

In 2019, what I considered to be a fight, became a nightmare. With the election of Brazil's former president, Jair Bolsonaro, the resources reduction became direct attacks: to Science and scientists, to Education and educators and even to Medicine and doctors. I saw most of the projects that I used to see during my childhood begin to ruin.

Despite all the problems I listed, and conciliating my studies with the fight for my rights, I successfully finished my High School and, in 2020 through SISU (another Lula's government project), got enrolled in a Computer Science course at the biggest Brazilian university, University of S\~ao Paulo. It was surely the consolidation of the beginning of my dream.

Since then, I keep researching, writing and studying a lot. I also have been talking a lot to my research colleagues and teachers. In common, their stories always state the same facts: the opportunities are much lesser than before. Outside my university, my old friends are researching but, at sometimes, fighting to keep their universities alive.

In the end of 2022, after the biggest Brazil's election, Lula was elected again, with the ``union and reconstruction'' slogan, proposing deep social transformation projects again. I clearly know that it doesn't mean instant solutions and improvements, but the end of the nightmare with changes. Changes like what I saw during my childhood. Changes that made my dream of becoming a scientist being a Brazilian become a dream and also a hope.
